\chapter*{CONCLUSION GÉNÉRALE}
\addcontentsline{toc}{chapter}{Conclusion générale}
\label{sec:conclusion}

Au terme de ce travail consacré à la conception et à la réalisation d'un système d'information intelligent au service de la valorisation de l'artisanat, il est possible de revenir sur le chemin parcouru et sur ce qu'il établit.

Le point de départ était un constat de terrain, non une intuition technologique. Le Village Artisanal Régional de Bafoussam dispose déjà d'un système d'information : un registre manuscrit, deux classeurs de suivi et la mémoire de son personnel. Ce système fonctionne, en ce sens qu'il permet à la structure de conduire son activité quotidienne. Il ne permet pas, en revanche, de rendre compte : ni de dire à un artisan ce qui lui revient lorsque le registre ne porte pas son nom, ni d'établir un taux de recouvrement sans un travail de rapprochement que personne n'entreprend en dehors d'un relevé annuel.

L'analyse conduite sur l'exercice~2026 a chiffré cette limite. Sur les 195~ventes de la période, 41 --- soit 21,0~\% --- portent un nom que rien ne rattache au parc locatif, ce qui représente 159~700~francs~CFA, un cinquième du chiffre d'affaires. La dette envers les artisans atteint 226~850~francs~CFA, soit 28,7~\% du chiffre d'affaires. Le taux de recouvrement des redevances s'établit à 14,73~\%, laissant 1~959~000~francs~CFA à recouvrer sur une imputation de 2~478~000~francs~CFA. Ces chiffres n'étaient pas connus de la structure avant ce travail ; ils constituent, indépendamment de toute solution logicielle, un résultat en eux-mêmes.

Un dernier constat dépasse les précédents en portée. La situation de caisse fait apparaître une différence de 162~525~francs~CFA entre les avoirs dus et les sommes détenues, à côté de laquelle figurent quatre sorties datées totalisant 120~000~francs~CFA. Il est vraisemblable que ces sorties expliquent l'essentiel de la différence ; il est certain que les supports actuels ne permettent pas de l'établir. Ce n'est pas le constat d'une irrégularité, c'est celui d'une question à laquelle la structure ne peut pas répondre. Un village artisanal qui manipule des recettes publiques doit pouvoir y répondre, et c'est ce que le journal de caisse immuable et l'arrêté journalier lui donnent.

La solution construite en réponse est une application modulaire de six modules à dépendance descendante, couvrant la chaîne complète de l'activité : de l'attribution d'un espace locatif jusqu'au reversement des sommes dues, en passant par le dépôt, la vente, l'écriture en caisse et la production des indicateurs. Trois propriétés en fondent la crédibilité dans une structure qui manipule des recettes publiques : les écritures financières y sont immuables et se corrigent par contre-passation ; les soldes y sont calculés et jamais saisis ; et toute action sensible y conserve l'identité de son auteur, constatée depuis le compte connecté et jamais choisie dans une liste.

Le volet d'intelligence artificielle prend la forme d'un assistant d'interrogation en langage naturel. Sa conception repose sur une décision antérieure à tout choix technique : une question portant sur un montant et une question portant sur un contenu ne sont pas traitées par le même mécanisme, et les deux mécanismes ne communiquent pas. Cette séparation, doublée d'un contrôle qui refuse tout chiffre absent des sources, garantit qu'aucun montant ne peut être produit par proximité textuelle. Mesuré sur 48~questions, le dispositif classe correctement 100~\% des questions, atteint un rappel@5 de 70~\% et refuse correctement 100~\% des questions sans réponse au corpus.

C'est sur ce dernier volet que le travail apporte son résultat le moins attendu. L'hypothèse selon laquelle une recherche vectorielle dense surpasserait une recherche par pondération de termes a été construite, mesurée, et infirmée : la branche dense perd dix points de rappel et fait tomber le taux de refus correct de 100~\% à 0~\%, parce que le modèle de plongement mobilisé est massivement anglophone tandis que les fiches du village font deux ou trois mots de français. Dans un espace où tout est également proche de tout, aucun seuil ne sépare rien. Ce résultat négatif, documenté et reproductible, vaut mieux qu'une fonctionnalité affirmée sans preuve : la première se démontre, la seconde s'énonce.

Il serait toutefois excessif de présenter cette solution comme achevée. Cinq limites en bornent la portée, dont trois appellent une suite : les prix ne représentent pas les conditionnements multiples ; l'espace artisan demeure au stade de la maquette, ce qui prive l'artisan de l'alerte de rupture qui lui est pourtant destinée ; et la couche d'autorisation repose sur un contrôle par action, sans politique déclarée au niveau des modèles. Ces limites sont consignées avec leur motif, de sorte qu'un repreneur sache ce qui a été délibérément laissé de côté et pourquoi.

Au-delà de ces points, la question décisive n'est pas technique. Un projet numérique antérieur, conduit dans la même structure et techniquement achevé, n'a jamais été déployé faute de budget d'hébergement. Ce précédent a commandé les choix de ce travail plus sûrement qu'aucune considération de performance : le cœur du système fonctionne sur un poste du village, en réseau local, sans abonnement ni connexion permanente, et seule la mise en ligne du portail engage une dépense annuelle, isolée et chiffrée pour que la coordination puisse en décider séparément.

Les perspectives ouvertes suivent trois directions. La première est fonctionnelle : la représentation des conditionnements, l'espace artisan, l'échéancier des redevances et les canaux de notification constituent une version~2 dont le périmètre est déjà documenté. La deuxième est méthodologique : le protocole d'évaluation construit ici est transposable à d'autres dispositifs d'interrogation en langage naturel, et l'expérience suggère qu'un modèle de plongement adapté au français mériterait d'être mesuré avant d'écarter définitivement la voie dense. La troisième est institutionnelle : le modèle de données ne présuppose pas l'unicité du village, et sa transposition aux autres villages artisanaux régionaux ne demanderait pas de refonte.

Un système d'information ne vaut, en définitive, que par la reprise qu'en fait la structure qui l'accueille. C'est le risque le plus élevé qui ait été identifié, et celui qu'aucune ligne de code ne traite. Le manuel utilisateur, la formation des agents et la désignation d'un référent interne y répondent partiellement. Le reste appartient au Village.
